Let
\[
G=\int_{-1}^1K(u)S(u)S(u)^\top\,du.
\]
The matrix \(G\) is positive definite: if \(v^\top Gv=0\), then the polynomial \(v^\top S(u)\) vanishes on a neighbourhood on which \(K\) is bounded below, and hence is identically zero, so \(v=0\). For \(h\le\min(t_0,1-t_0)\), the change of variables \(a=t_0+hu\) and Lemma~\texttt{lem:marginal-treatment-positivity} give
\[
D_h=
\int_{-1}^1K(u)S(u)S(u)^\top f_A(t_0+hu)\,du
\succeq cG.
\]
Thus \(D_h\) is invertible and \(\|D_h^{-1}\|_{\mathrm{op}}\le\{c\lambda_{\min}(G)\}^{-1}\). If \(q\) has degree at most \(p\), write \(q(t_0+hu)=\gamma_q^\top S(u)\); its intercept is \((\gamma_q)_0=q(t_0)\). It follows algebraically that
\[
\int W_h(a)q(a)f_A(a)\,da
=e_0^\top D_h^{-1}D_h\gamma_q
=q(t_0).
\]
By Lemma~\texttt{lem:oracle-pseudo-outcome-identity},
\[
\vartheta_h=
\int W_h(a)\theta_P(a)f_A(a)\,da.
\]
Let \(T_{t_0}\) be the Taylor polynomial of \(\theta_P\) at \(t_0\), of degree \(p\) when \(\alpha\notin\mathbb N\), and degree \(p-1\) when \(\alpha=p\in\mathbb N\). The latter polynomial is also reproduced because reproduction holds through degree \(p\). Lemma~\texttt{lem:partial-mean-holder-transfer} and Taylor's formula under the frozen seminorm convention give
\[
|\theta_P(a)-T_{t_0}(a)|
\le C|a-t_0|^\alpha.
\]
Moreover,
\[
\int|W_h(a)|f_A(a)\,da
\le
\|e_0^\top D_h^{-1}\|
\int_{-1}^1K(u)\|S(u)\|f_A(t_0+hu)\,du
\le C.
\]
Polynomial reproduction therefore yields
\[
|\vartheta_h-\theta_P(t_0)|
\le Ch^\alpha.
\]
Only the two-sided bounds on \(f_A\) were used; no smoothness of the marginal treatment density entered.